% BHU PYQ -- Manually transcribed & error-corrected
% Paper: BCH-321: Financial Analysis
% Exam: B.Com. (Hons.) Semester VI Examination, 2015-16

\documentclass[11pt,a4paper]{article}
\usepackage[a4paper,top=2.5cm,bottom=2.5cm,left=2.8cm,right=2.8cm,headheight=14pt]{geometry}
\usepackage{amsmath,amssymb}
\usepackage[T1]{fontenc}
\usepackage[utf8]{inputenc}
\usepackage{lmodern,microtype,fancyhdr,enumitem,hyperref,lastpage,parskip}

\hypersetup{
    colorlinks=true,
    urlcolor=black,
    linkcolor=black,
    pdftitle={BCH-321: Financial Analysis -- BHU B.Com. (Hons.) Sem VI 2015-16}
}

\pagestyle{fancy}
\fancyhf{}
\renewcommand{\headrulewidth}{0.4pt}
\renewcommand{\footrulewidth}{0.4pt}
\fancyhead[L]{\small\textit{Banaras Hindu University}}
\fancyhead[R]{\small\textit{BCH-321: Financial Analysis}}
\fancyfoot[C]{\small Page \thepage\ of \pageref{LastPage}}

\newlist{parts}{enumerate}{2}
\setlist[parts,1]{label=(\alph*),leftmargin=*,itemsep=5pt,topsep=3pt}
\setlist[parts,2]{label=(\roman*),leftmargin=*,itemsep=3pt,topsep=2pt}
\newcommand{\pts}[1]{\hfill\textit{[#1]}}

\begin{document}

\begin{center}
  {\large\bfseries BANARAS HINDU UNIVERSITY}\\[2pt]
  {\normalsize B.Com. (Hons.) --- Semester VI Examination, 2015-16}\\[6pt]
  \rule{0.8\linewidth}{0.5pt}\\[6pt]
  {\large\bfseries Paper No. BCH-321: Financial Analysis}\\[4pt]
  {\normalsize Time: 3 Hours\hfill Maximum Marks: 70}
\end{center}

\rule{\linewidth}{0.4pt}

\smallskip
\noindent\textit{Instructions: Answer all questions. All questions carry equal marks (14 marks each).}

\bigskip

\noindent\textbf{Question 1.} \\
Explain the meaning of ``Analysis of Financial Statements''. Describe the different tools of financial analysis.\pts{14}

\centerline{\textbf{OR}}

What are the objectives of Financial Analysis? State the stakeholders that are interested in the analysis of financial statements.\pts{14}

\medskip\rule{\linewidth}{0.2pt}

\noindent\textbf{Question 2.} \\
Discuss the significance of ratio analysis. Also throw light on its limitations.\pts{14}

\centerline{\textbf{OR}}

\begin{parts}
  \item The following is the Balance Sheet of New India Ltd. for the year ending 31st December, 2014:\pts{8}
  
  \begin{center}
  \small
  \renewcommand{\arraystretch}{1.2}
  \begin{tabular}{|l|r|l|r|}
  \hline
  \textbf{Liabilities} & \textbf{Amount (Rs.)} & \textbf{Assets} & \textbf{Amount (Rs.)} \\ \hline
  9\% Preference Share Capital & 5,00,000 & Goodwill & 1,00,000 \\
  Equity Share Capital & 10,00,000 & Land \& Building & 6,50,000 \\
  8\% Debentures & 2,00,000 & Plant & 8,00,000 \\
  Long-term Loan & 1,00,000 & Furniture & 1,50,000 \\
  Bills Payable & 60,000 & Bills Receivables & 70,000 \\
  Creditors & 70,000 & Debtors & 90,000 \\
  Bank Overdraft & 30,000 & Bank Balance & 45,000 \\
  Outstanding Expenses & 5,000 & Short-term Investments & 25,000 \\
  & & Pre-paid Expenses & 5,000 \\
  & & Stock & 30,000 \\ \hline
  \textbf{Total} & \textbf{19,65,000} & \textbf{Total} & \textbf{19,65,000} \\ \hline
  \end{tabular}
  \end{center}

  From the above Balance Sheet, calculate and comment on the following Ratios:
  \begin{parts}
    \item Current Ratio
    \item Acid-test Ratio
    \item Liquidity Ratio
  \end{parts}
  
  \item The capital of ABC Ltd. is as follows:\pts{6}
  \begin{center}
  \renewcommand{\arraystretch}{1.1}
  \begin{tabular}{lr}
    9\% Preference Shares of Rs. 10 each & Rs. 3,00,000 \\
    Equity Shares of Rs. 10 each & Rs. 8,00,000 \\ \hline
    \textbf{Total} & \textbf{Rs. 11,00,000} \\ \hline
  \end{tabular}
  \end{center}
  
  The accountant has ascertained the following information:
  \begin{itemize}[itemsep=2pt,topsep=2pt]
    \item Profit after tax: Rs. 2,70,000 (tax rate is 60\%)
    \item Equity Dividend Paid: 20\% (Rs. 1,60,000)
    \item Market Price Per Equity Share: Rs. 40
  \end{itemize}
  You are required to calculate the following:
  \begin{parts}
    \item Earning Per Share (EPS)
    \item Price Earning Ratio (P/E Ratio)
  \end{parts}
\end{parts}

\medskip\rule{\linewidth}{0.2pt}

\pagebreak

\noindent\textbf{Question 3.} \\
Give a detailed format of a cash flow statement. How does it differ from a funds flow statement?\pts{14}

\centerline{\textbf{OR}}

The Balance Sheet of Paschim Corporation Ltd. as on 31st December, 2014 and 2015 is as follows:\pts{14}

\begin{center}
\small
\renewcommand{\arraystretch}{1.2}
\begin{tabular}{|l|c|c|l|c|c|}
\hline
\textbf{Liabilities} & \textbf{2014 (Rs.)} & \textbf{2015 (Rs.)} & \textbf{Assets} & \textbf{2014 (Rs.)} & \textbf{2015 (Rs.)} \\ \hline
11\% Cum. Pref. Shares & -- & 30,000 & Land \& Building & 60,000 & 50,000 \\
Equity Shares & 1,10,000 & 1,20,000 & Plant & 30,000 & 50,000 \\
Reserve & 4,000 & 4,000 & Sundry Debtors & 40,000 & 48,000 \\
Profit \& Loss A/c & 2,000 & 2,400 & Stock & 60,000 & 70,000 \\
9\% Debentures & 12,000 & 14,000 & Bank & 2,400 & 7,000 \\
Provision for Taxation & 6,000 & 8,400 & Cash & 600 & 1,000 \\
Proposed Dividend & 10,000 & 11,600 & & & \\
Current Liabilities & 49,000 & 35,600 & & & \\ \hline
\textbf{Total} & \textbf{1,93,000} & \textbf{2,26,000} & \textbf{Total} & \textbf{1,93,000} & \textbf{2,26,000} \\ \hline
\end{tabular}
\end{center}

You are required to prepare a schedule of changes in working capital and a statement of Flow of Funds.

\medskip\rule{\linewidth}{0.2pt}

\noindent\textbf{Question 4.} \\
Explain the meaning and need of financial forecasting in business.\pts{14}

\centerline{\textbf{OR}}

What is a Projected Income Statement? Explain the steps involved in preparing a projected income statement.\pts{14}

\medskip\rule{\linewidth}{0.2pt}

\noindent\textbf{Question 5.} \\
Discuss the various methods suggested for determining the valuation of shares.\pts{14}

\centerline{\textbf{OR}}

Compute the goodwill of a company from the following information:\pts{14}
\begin{enumerate}[label=(\roman*),itemsep=2pt,topsep=2pt]
  \item Capital employed: Rs. 3,00,000
  \item Outside liabilities: Rs. 90,000
  \item Profits of the preceding three years: Rs. 30,000, Rs. 28,000, and Rs. 38,000
  \item The normal rate of return: 10\%
  \item Interest earned Rs. 3,000 p.a. has been included in the profit of the last three years
  \item The fair remuneration to the proprietor of the business may be assumed as Rs. 800 p.m.
\end{enumerate}

Calculate goodwill by:
\begin{parts}
  \item Capitalisation of average profit method
  \item Capitalisation of super profit method
\end{parts}

\vfill
\rule{\linewidth}{0.4pt}
\begin{center}\small
\href{https://bhu-exam.vercel.app/}{Click here to visit our website}\\[4pt]
\href{https://me-aryan.vercel.app/}{Click here to visit our contributor}\\[4pt]
\href{https://quashan-me.vercel.app/}{Click here to visit our contributor}
\end{center}

\end{document}
